\chapter{Related Literature}
\label{Chapter2}

\section{Classical TWFE approach and its problems}

Recent work on DiD begins from a now well-established point. With multiple periods, staggered adoption, and heterogeneous treatment effects, the \emph{two-way fixed effects} (TWFE) estimator no longer admits the canonical \(2\times 2\) interpretation. In its most common form,
\begin{equation}
Y_{it} = \alpha_i + \lambda_t + \beta D_{it} + \varepsilon_{it},
\label{eq:twfe_static}
\end{equation}
where \(Y_{it}\) is the outcome of interest, \(\alpha_i\) and \(\lambda_t\) are unit and period fixed effects, and \(D_{it}\) indicates treatment. In the \(2\times 2\) case, under parallel trends and no anticipation, \(\beta\) coincides with the usual DiD estimand. In staggered settings, that interpretation breaks down. \textcite{roth2023s} synthesize the theoretical literature and argue that the main insight is not only about weighting. It also concerns which comparisons identify the effect, how dynamic specifications should be interpreted, and how covariates enter the design. The exposition below follows that three-part division.

\subsection{Weighting and the Goodman-Bacon decomposition}

\textcite{goodman2021difference} shows that, in the canonical staggered-adoption TWFE design without covariates, the coefficient can be written as a weighted average of admissible \(2\times 2\) DiD estimators:
\begin{equation}
\hat{\beta}^{TWFE}
=
\sum_{j} \omega_j \hat{\beta}^{2\times 2}_j,
\qquad
\sum_j \omega_j = 1.
\label{eq:bacon_generic_short}
\end{equation}
The main implication is that TWFE combines multiple cohort-time comparisons rather than recovering a single ATT-type parameter.\footnote{In \textcite{goodman2021difference}, the relevant comparisons include treated-versus-never-treated and earlier-treated-versus-later-treated \(2\times 2\) designs. The notation \(\omega_j\) in equation~\eqref{eq:bacon_generic_short} is used here as a compact generic label for the Goodman-Bacon weights. A related insight appears in \textcite{dechaisemartin2020twfe}, who study a broader two-way fixed-effects setting and show that the TWFE estimand can be expressed as a weighted sum of group-period treatment effects with weights that sum to one but may be negative.} Once treatment effects vary across cohorts or over exposure time, some of those comparisons become difficult to interpret as clean causal contrasts. The coefficient may then depart from the effect researchers intend to estimate, sometimes by a great deal. The applied consequence is that a single TWFE coefficient can conceal considerable internal heterogeneity across the $2\times 2$ contrasts that construct it.

\subsection{Dynamic specifications and event-study contamination}

A parallel problem arises in dynamic \emph{event-study} specifications:
\begin{equation}
Y_{it}
=
\alpha_i + \lambda_t
+ \sum_{\ell \neq -1} \mu_{\ell}\mathbf{1}\{t-E_i=\ell\}
+ u_{it},
\label{eq:twfe_event_study_short}
\end{equation}
where \(E_i\) denotes the treatment adoption period and \(\ell=t-E_i\) is event time. A common interpretation is that \(\mu_\ell\) measures the average treatment effect at horizon \(\ell\). \textcite{sun2021estimating} show that this interpretation generally breaks down under staggered adoption with heterogeneous treatment effects across cohorts: the coefficient for horizon \(\ell\) generally loads not only on the effect at that horizon but also on effects from other relative periods. Their interaction-weighted estimator removes this contamination by averaging cohort-specific effects only across cohorts that identify the same horizon. This matters because applied work routinely uses lead coefficients to assess pre-trends and lag coefficients to describe timing and persistence. Under treatment-effect heterogeneity, those coefficients may misstate both the credibility of the design and the shape of the treatment profile they are supposed to recover.

\subsection{Modern DiD estimators: identification before aggregation}

Recent work has reoriented the discussion toward designs that identify treatment effects at the group-time level \emph{before} any aggregation. In \textcite{callaway2021difference}, the basic object is the group-time average treatment effect,
\begin{equation}
ATT(g,t)
=
\mathbb{E}\!\left[ Y_t(g)-Y_t(0)\mid G_g=1 \right],
\label{eq:att_gt_short}
\end{equation}
where \(G_g\) denotes the group first treated in period \(g\). The key contribution lies in the identification strategy: treatment effects are first defined for each \((g,t)\) cell and only afterward aggregated into summary parameters. This makes the comparison group explicit. In practice, the choice appears in the distinction between \emph{never-treated} and \emph{not-yet-treated} controls, which correspond to different versions of the parallel trends assumption. \textcite{marcus2021role} formalize the trade-off between these two comparison groups in event-study settings: the not-yet-treated group is typically larger and more efficient but relies on a stronger parallel-trends assumption (pre-trends must hold for later cohorts in their untreated years), while the never-treated group imposes a weaker assumption but at the cost of efficiency when the untreated pool is thin. \textcite{callaway2021difference} implements both versions, leaving the choice to the applied researcher as a substantive design decision rather than a technical default.

Complementary estimators build on the same principle of identifying $ATT(g,t)$ first. \textcite{borusyak2024revisiting} propose an imputation approach that explicitly models untreated outcomes and infers treatment effects by extrapolation, which \textcite{gardner2022two} implements as a two-stage estimator. \textcite{wooldridge2021two,wooldridge2025two} shows that a saturated version of TWFE --- his Extended Two-Way Fixed Effects (ETWFE) --- recovers the $ATT(g,t)$ targets of Callaway--Sant'Anna through a standard regression with interaction terms. \textcite{arkhangelsky2021synthetic} introduce Synthetic Difference-in-Differences, which combines synthetic-control re-weighting with DiD, and has since become a further alternative. We do not use Synthetic DiD in our reanalysis, but mention it here because it forms part of the modern methodological landscape the applied researcher now faces. The common thread across these estimators is that they make explicit which comparisons identify the effect and at which level aggregation occurs; they differ in how they construct the counterfactual and in what assumptions identification requires.

\subsection{Parallel trends: sensitivity, not pre-testing}

Even with more interpretable estimators, identification still depends on some version of parallel trends. Applied work often evaluates that assumption through pre-trend tests in dynamic specifications. \textcite{roth2022pretest} shows why this practice is weaker than it often appears: failure to reject pre-trends may simply reflect low power, and conditioning the analysis on passing that test induces pre-test bias. \textcite{ghanem2024selection} further clarify when parallel trends is itself defensible, linking the assumption to selection models and showing that it holds only under specific restrictions on unobserved heterogeneity.

\textcite{rambachan2023more} propose a different approach, replacing the pass-or-fail logic of pre-trend tests with explicit sensitivity analysis. Rather than treat exact parallel trends as something established by insignificant leads, they ask how conclusions change under bounded deviations from that benchmark. This shift is important for the broader literature because it moves credibility assessment away from mechanical hypothesis testing and toward transparent robustness to plausible violations. In the applied reanalysis below, we implement the Rambachan--Roth sensitivity framework on every article that has a dynamic event study, alongside the estimator-swap exercise.

\subsection{Covariates as part of the design}

Covariate choices are also part of identification. In modern DiD designs, adding controls replaces an unconditional parallel trends assumption with a conditional one, so the design begins to depend on which covariates enter and how they enter. In \textcite{callaway2021difference}, this conditioning is built around pre-treatment covariates, and \(ATT(g,t)\) can be implemented through outcome regression, inverse probability weighting, or doubly robust estimands --- the last of which is the \textcite{sant2020doubly} estimator whose efficiency gains hold even when parallel trends is unconditional. This is conceptually different from the familiar practice of adding controls to a TWFE regression and interpreting the result as a more flexible version of the same estimand.

The distinction becomes sharper once one considers time-varying covariates. As emphasized by \textcite{caetano2022difference}, if treatment can affect these covariates, conditioning on their realized post-treatment values may amount to conditioning on mediators or selection variables rather than on untreated characteristics. Even when treatment does not affect them directly, \textcite{caetano_callaway2024_hiddenlinearity} and \textcite{wooldridge2025two} show that standard TWFE implementations can still impose additional and often opaque restrictions. After removing unit fixed effects, the regression may condition only on changes in time-varying covariates, while time-invariant covariates disappear from the transformed specification altogether. The lesson from this literature is that conditioning choices alter the maintained assumptions behind the design and can affect interpretation in ways that are easy to miss in routine empirical practice.

\subsection{Implementation as a finite-sample concern}

Implementation details can also matter in finite samples. In principle, differences across modern DiD estimators should reflect differences in identification, aggregation, or target parameters. In practice, numerical choices such as balancing rules, support restrictions, and first-step estimation can also affect results, especially in highly unbalanced panels or designs with weak overlap. This point belongs in the literature review because it clarifies that disagreements across estimates need not always reflect purely conceptual differences: some arise because distinct commands operationalize the same estimand differently in finite samples. Primary evidence for this claim appears in Section~\ref{sec:li_sensitivity} (Lesson~10), which documents three balancing rules under a nominally identical CS-DID specification producing qualitatively different post-treatment profiles in one of the reanalyzed studies.

\section{The reanalysis movement in economics and adjacent fields}

This paper sits inside a broader disciplinary movement to re-examine published findings using more credible tools. The psychology replication crisis that crystallized with the \textcite{openscience2015reproducibility} and \textcite{camerer2018evaluating} reproducibility projects has reshaped how applied economists think about the empirical record: credibility of a finding is assessed not only at the point of first publication, but over a longer horizon that includes replication, re-estimation under alternative methods, and generalization to adjacent settings \parencite{vivalt2020specification}. The present exercise inherits that framing. The contribution of this paper is not to argue for the credibility of any individual DiD estimate; it is to ask how the modal DiD estimate in a defined applied literature moves when the estimator that originally produced it is replaced.

Within economics, three recent reanalysis studies define the immediate predecessors of this exercise. Rather than review them chronologically, we organize them by the distinct contribution each makes to the present paper.

\textbf{\textcite{baker2022much}} --- \textit{applied consequence demonstration.} BLW connect the theoretical critique of staggered TWFE to applications in corporate finance and accounting. Their contribution is less about the scale of the exercise than about establishing, in a defined applied literature, that replacing TWFE with heterogeneity-robust estimators can change magnitudes, dynamic profiles, and specification choices such as endpoint binning. The value for us is methodological: BLW shows that the estimator-swap exercise is a productive unit of analysis, and their diagnostic framework for staggered-adoption finance designs transfers directly to our setting.

\textbf{\textcite{chiu2023and}} --- \textit{identification-as-fragility insight.} Chiu, Lau, Lau, and Xu reanalyze 49 studies published in the \emph{American Political Science Review}, \emph{American Journal of Political Science}, and \emph{The Journal of Politics} that use TWFE for causal inference in observational panel settings. Their contribution reshapes the reanalysis agenda. In many applications, heterogeneity-robust estimators remain qualitatively close to TWFE; some show sign reversals or precision losses, but those are not the modal outcome. The more consequential finding concerns identification itself: only a subset of studies produces direct evidence against parallel trends, but many more are too underpowered to rule out violations large enough to alter the practical meaning of the estimate. Chiu et al.\ thereby shift the emphasis of reanalysis from \emph{weighting correction} to \emph{identification credibility} --- a shift we inherit and extend via the HonestDiD sensitivity analysis in Lesson~5.

\textbf{\textcite{brodeur2026reproducibility}} --- \textit{robustness-vs-reproducibility distinction.} Brodeur and co-authors reanalyze 110 recent papers published in top economics and political-science journals under mandatory data and code sharing policies. Their contribution is taxonomic: they distinguish \emph{computational reproducibility} (can the published numbers be recovered from the supplied code?) from \emph{robustness} (do the substantive conclusions survive alternative but reasonable analytical choices?). Reproducibility rates are high, but robustness rates are substantially lower, and the gap is a diagnostic about where the literature is most fragile. For the present paper, that distinction is clarifying: our exercise is a robustness exercise, not a reproducibility one. We do not primarily test whether each paper's original code recovers its published numbers; we test whether its conclusions survive re-estimation with estimators developed after its publication.

\subsection{The gap this paper fills}
\label{sec:ch2_gap}

\textcite{baker2022much} covers corporate finance and accounting. \textcite{chiu2023and} covers political science. \textcite{brodeur2026reproducibility} covers cross-field robustness. \textbf{None of them, and no comparable study we are aware of, examines health economics at scale or systematically documents the contribution of specific design margins --- binning, geographic conditioning, dynamic aggregation, covariate handling, and software defaults --- to the apparent divergence between TWFE and modern estimators in any single applied subfield.} This is the gap the paper fills.

Health economics is a natural site for this exercise for three reasons. First, DiD is the workhorse empirical design in the subfield and its use has accelerated over the past decade \parencite{goldsmith2024tracking}. Second, the policy stakes are unusually explicit: many of the re-estimated articles inform concrete welfare or cost-effectiveness calculations in which the magnitude of the treatment effect, not only its sign, determines whether the policy clears the adoption threshold (see Section~\ref{Chapter4}, Lesson~1). Third, the health-econ literature has developed its own applied best-practice guidance for DiD \parencite{wing2018designing}, which means there is an established methodological anchor against which our reanalysis findings can be read. The three predecessors named above do not provide that anchor for health.

This paper therefore trades breadth across fields for depth across design margins within a single field. Rather than replicate the strategy of examining estimator swaps only, we document \emph{how} the swap changes the estimate --- which margin of the design carries the change, and whether the change reflects genuine correction, estimand redefinition, or implementation artefact. The ten lessons in Section~\ref{Chapter4} are organized around those margins.
